\documentclass{article}
\usepackage[T1]{fontenc}
\usepackage{lmodern}
\usepackage{graphicx}
\usepackage{amsmath,amssymb}
\usepackage[a4paper,margin=2cm]{geometry}
\usepackage[absolute,overlay]{textpos}
\setlength{\TPHorizModule}{1mm}
\setlength{\TPVertModule}{1mm}
\usepackage[indonesian]{babel}
\usepackage{hyperref}
\setlength{\emergencystretch}{2em}
% unit: Halaman 15

\begin{document}

% === Halaman 15 (llm) ===

\begin{flushright}
\fbox{Enkripsi: $c = m^e \pmod{n}$}
\end{flushright}

\begin{itemize}
    \item Bob mengenkripsi setiap blok plainteks dengan menggunakan kunci publik Alice ($e = 79$, $n = 3337$):
    \begin{align*}
        m_1 = 0704 & \to c_1 = 704^{79} \pmod{3337} = 328; \\
        m_2 = 1111 & \to c_2 = 1111^{79} \pmod{3337} = 301; \\
        m_3 = 1400 & \to c_3 = 1400^{79} \pmod{3337} = 2653; \\
        m_4 = 1108 & \to c_4 = 1108^{79} \pmod{3337} = 2986; \\
        m_1 = 204 & \to c_5 = 204^{79} \pmod{3337} = 1164;
    \end{align*}

    Cipherteks: $C = 0328, 0301, 2653, 2986, 1164$

    \item Bob mengirim cipherteks $C$ kepada Alice
\end{itemize}

\end{document}
